% Generated draft: requires mathematical and source review.
\documentclass[11pt]{article}
\usepackage[margin=25mm]{geometry}
\usepackage{fontspec}
\setmainfont{Latin Modern Roman}
\usepackage{amsmath,amssymb,mathtools,unicode-math}
\setmathfont{Latin Modern Math}
\usepackage{enumitem,needspace,xurl,hyperref}
\setlength{\parindent}{0pt}
\setlength{\parskip}{5pt}
\setlength{\emergencystretch}{3em}
\newcommand{\N}{\mathbb{N}}
\newcommand{\Z}{\mathbb{Z}}
\newcommand{\Q}{\mathbb{Q}}
\newcommand{\R}{\mathbb{R}}
\newcommand{\C}{\mathbb{C}}
\newcommand{\F}{\mathbb{F}}
\newcommand{\A}{\mathbb{A}}
\newcommand{\PP}{\mathbb{P}}
\newcommand{\E}{\mathbb{E}}
\newcommand{\eps}{\varepsilon}
\newcommand{\dd}{\,\mathrm{d}}
\begin{document}
\section*{UnsolvedMath ID: 30001476}
\textbf{Source problem code:} catalogue entry
\subsection{Definitions and mathematical statement}
Let \(\F_2=\Z/2\Z\). For a compact flat three-manifold \(M\), let \(H_2(M;\F_2)\) denote its second homology group with mod-two coefficients. Given a specified class \(\alpha\in H_2(M;\F_2)\) and a specified collection \(\mathcal C(M,\alpha)\) of admissible surface representatives of \(\alpha\), define
\[
 m(M,\alpha)=\inf\{\operatorname{Area}(T):T\in\mathcal C(M,\alpha)\}.
\]
An admissible representative \(S\) is an area minimizer in \(\alpha\) if
\[
 [S]=\alpha,\qquad \operatorname{Area}(S)=m(M,\alpha).
\]
The supplied record does not specify the admissible category, for example whether competitors are smooth embedded surfaces or mod-two integral cycles. Thus \(\mathcal C(M,\alpha)\) remains to be fixed from the source context.

Write \(M_X(h)\) for the compact flat three-manifold called the \(X\) Cosm in the supplied record, where
\[
 X\in\{\mathrm{Tri},\mathrm{Tetra},\mathrm{Hexa}\},\qquad h\in I_X,
\]
and \(h\) is its height parameter. The parameter domains \(I_X\), the quotient constructions, and the metric normalizations are not supplied. Let \(P_X(h)\) denote the horizontal planar surface referred to in the record, and let \(S_X(h)\) denote the corresponding catenoid-like candidate, with names assigned by
\[
 S_{\mathrm{Tri}}=\mathrm{rPD},\qquad
 S_{\mathrm{Tetra}}=\mathrm{tP},\qquad
 S_{\mathrm{Hexa}}=\mathrm{H}.
\]
These names identify source candidates; the record does not provide their constructions or their dependence on \(h\).

\textbf{Question, with unresolved source parameters.} For each of these three families, is there a division of the admissible height range into a planar range \(I_X^{\mathrm{plan}}\) and a complementary catenoid-like range \(I_X^{\mathrm{cat}}\) such that, in the intended homology class \(\alpha_X(h)\),
\[
 \begin{aligned}
 \operatorname{Area}(P_X(h))&=m(M_X(h),\alpha_X(h))
 &&\text{for every }h\in I_X^{\mathrm{plan}},\\
 \operatorname{Area}(S_X(h))&=m(M_X(h),\alpha_X(h))
 &&\text{for every }h\in I_X^{\mathrm{cat}}?
 \end{aligned}
\]
The intended representatives must belong to the indicated class wherever they are compared. Their homology classes, the height ranges, and the treatment of any transition value are not specified in the supplied record; in particular, no assertion of uniqueness or of strict inequalities between competing areas is included here.

For the Tetra Cosm, the record additionally refers to two remaining nonzero homology classes. Denote the class intended for \(\mathrm{CLP}\) by \(\beta(h)\), and that intended for \(\mathrm{tD}\) by \(\gamma(h)\). Are the named minimal surfaces area minimizers in those respective classes, that is,
\[
 \operatorname{Area}(\mathrm{CLP}(h))=m(M_{\mathrm{Tetra}}(h),\beta(h)),
 \qquad
 \operatorname{Area}(\mathrm{tD}(h))=m(M_{\mathrm{Tetra}}(h),\gamma(h))?
\]
Here the notation records the intended dependence on the ambient metric without asserting a supplied construction. The parameter values for which these assertions are intended, the identification of the two classes, and the realizations of the surfaces remain unspecified. Consequently, these formulas give the variational content of the supplied questions, but a fully determined formulation requires the missing geometric and parameter data.
\subsection{Short English statement}
The question asks whether varying the height of the Tri, Tetra, and Hexa Cosms changes the mod-two area minimizer from a horizontal planar surface to, respectively, \(\mathrm{rPD}\), \(\mathrm{tP}\), and \(\mathrm{H}\), and whether \(\mathrm{CLP}\) and \(\mathrm{tD}\) minimize area in the two other indicated nonzero classes of the Tetra Cosm. The supplied record leaves the metrics, parameter ranges, transition behavior, surface constructions, and competitor category unresolved.
\subsection{Sources}
\begin{itemize}
\item Local UnsolvedMath registry: supplied complete source record, permanent ID \(30001476\), exact source problem code \texttt{catalogue entry}, title ``Homological Area Minimizers in Compact Flat Three-Manifolds,'' in the Oberwolfach Reports Open Problems collection. The record gives proposed year 2010 and identifies no individual proposer.
\item Supplied report citation: \emph{Progress in Surface Theory} (2010), Oberwolfach Reports, pp.\ 1225--1318, DOI \texttt{10.4171/owr/2010/21}. \url{https://doi.org/10.4171/owr/2010/21}. This is also the supplied literature-source entry labeled ``OWR 2010/21.''
\item Additional supplied literature-source entry, labeled ``Tetrahedral cosmology geometry'': \url{https://doi.org/10.1051/jp1:1992276}. No further bibliographic details are supplied in the record.
\item Literature check date: \texttt{2026-08-21}. The supplied dated curation labels the problem open and states that the report presents the \(\mathrm{CLP}\) and \(\mathrm{tD}\) area-minimization assertions as conjectures, while the located geometric work does not prove the required minimizing property. This is supplied dated curation, not an independently verified literature assessment.
\item The registry labels the statement \texttt{corrected\_verified} and reports that its mathematical content was checked against the official Oberwolfach report, with notation, context, and question form normalized. That verification claim is supplied provenance; no independent source inspection or mathematical verification is asserted here.
\end{itemize}
\end{document}
